\section{Results}
\label{sec:results}

This section must be populated only from machine-readable experiment outputs. Numerical values should not be typed manually into the manuscript.

\subsection{Deterministic Forecasting}
\todo{Insert generated multi-dataset, multi-horizon comparison table. Discuss mean$\pm$std performance, not only the best single run. Report cases where TrustKAN does not rank first.}

\IfFileExists{tables/benchmark_summary.tex}{\input{tables/benchmark_summary.tex}}{}

\subsection{Probabilistic Forecasting and Calibration}
\todo{Insert static conformal and quantile-forecast results: empirical coverage, interval width, interval score, and calibration deviation. Compare static, rolling, and adaptive calibration.}

\subsection{Adaptive Conformal Performance under Drift}
\todo{Report rolling coverage trajectories and interval-width trade-offs under temporal drift. Quantify whether adaptive calibration improves coverage stability and whether it pays for this through excessive interval expansion.}

\begin{figure}[t]
\centering
\IfFileExists{figures/rolling_coverage.pdf}{\includegraphics[width=\columnwidth]{rolling_coverage.pdf}}{\fbox{\parbox{0.9\columnwidth}{\centering Rolling coverage figure will be generated from experiment outputs.}}}
\caption{Rolling empirical coverage under static, rolling, and adaptive conformal calibration.}
\label{fig:rollingcoverage}
\end{figure}

\subsection{Reliability--Error Association}
\todo{Report Spearman/Pearson association, reliability-bin errors, and AUROC/AUPRC for identifying high-error predictions. A reliability score that does not predict realized error should not be presented as a core contribution.}

\begin{figure}[t]
\centering
\IfFileExists{figures/reliability_vs_error.pdf}{\includegraphics[width=\columnwidth]{reliability_vs_error.pdf}}{\fbox{\parbox{0.9\columnwidth}{\centering Reliability--error figure will be generated from experiment outputs.}}}
\caption{Relationship between estimated reliability and realized forecasting error.}
\label{fig:reliabilityerror}
\end{figure}

\subsection{Selective Prediction}
\todo{Insert risk--coverage curves and quantify AURC, retained-set RMSE/MAE, abstention rate, and large-error removal rate. Compare against simple uncertainty-only and shift-only selection.}

\begin{figure}[t]
\centering
\IfFileExists{figures/risk_coverage.pdf}{\includegraphics[width=\columnwidth]{risk_coverage.pdf}}{\fbox{\parbox{0.9\columnwidth}{\centering Risk--coverage figure will be generated from experiment outputs.}}}
\caption{Risk--coverage behaviour for selective forecasting.}
\label{fig:riskcoverage}
\end{figure}

\subsection{Ablation Study}
\todo{Report the predefined A0--A9 ablations. Remove any claimed contribution whose ablation does not show incremental and reproducible value.}

\subsection{Interpretability and Explanation Stability}
\todo{Visualize representative KAN functions and quantify stability across seeds, periods, and perturbations. Distinguish physically plausible observations from unsupported causal interpretation.}

\subsection{Robustness, Extremes, and Efficiency}
\todo{Report noise, missing-data, extreme-event, and computational experiments. Include failure cases and performance/complexity trade-offs.}
